\documentclass[11pt,a4paper]{article}
\usepackage[utf8]{inputenc}
\usepackage[T1]{fontenc}
\usepackage[english]{babel}
\usepackage{geometry}
\usepackage{longtable}
\usepackage{array}
\usepackage{amsmath}
\usepackage{hyperref}
\geometry{margin=2.5cm}

\title{Technical Documentation of \texttt{hydro\_1D\_IS.f}}
\author{PICHYDRO-1D\_IS Project}
\date{\today}

\begin{document}
\maketitle
\tableofcontents
\newpage

\section{Scope and Objective}
This document describes, in detail, how the program implemented in
\texttt{src/hydro\_1D\_IS.f} works. The code solves relativistic ideal
hydrodynamics for heavy-ion collisions using a Particle-In-Cell (PIC)
strategy, with an event-by-event initial state and a freeze-out hypersurface
built with Cornelius.

The main objectives of this manual are:
\begin{itemize}
\item explain the global control flow of the program;
\item describe the physical variables and their numerical representation;
\item document input files, internal phases, and expected outputs;
\item provide a maintenance-oriented map of key subroutines.
\end{itemize}

\section{Program Architecture}
\subsection{Entry point}
The execution starts in \texttt{program hydro}. The program:
\begin{enumerate}
\item reads \texttt{../input/inhydro};
\item validates run options (current branch requires \texttt{kslab=8});
\item initializes nuclear/EoS context (\texttt{cnucl}, \texttt{read\_chiral\_eos});
\item calls \texttt{picrel}, which contains the full time evolution.
\end{enumerate}

\subsection{Memory layout and shared state}
The code is Fortran-77 style and relies heavily on \texttt{COMMON} blocks.
The most important storage design is:
\begin{itemize}
\item \texttt{/plcm1/, /plcm2/, /plcm3/}: large contiguous arrays holding field
      variables through \texttt{equivalence};
\item \texttt{/plcm4/}: marker particle coordinates;
\item \texttt{/pscm1/}: global simulation parameters (mesh, time step, counters,
      units, geometry);
\item \texttt{/newstf/, /option/, /numcmm/}: run options and flags;
\item \texttt{/prevent/, /posevent/}: pre/post-cycle fields used for hypersurface
      reconstruction.
\end{itemize}

\section{Physical and Numerical Variables}
The main hydrodynamic fields on the Eulerian mesh are:
\begin{itemize}
\item \texttt{capn}: baryon density-like quantity ($N^0$ in the lab frame);
\item \texttt{e}: energy density-like quantity ($T^{00}$ in the lab frame);
\item \texttt{p}: pressure;
\item \texttt{xmom, ymom, zmom}: momentum densities ($T^{0x},T^{0y},T^{0z}$);
\item \texttt{u, v, w}: flow velocity components.
\end{itemize}

Marker particles are represented by:
\begin{itemize}
\item \texttt{px, py, pz}: continuous particle coordinates in mesh units.
\end{itemize}

\subsection{Units}
The source contains multiple comments from different development stages.
In the current branch, the active event-by-event path (\texttt{kslab=8})
works with energies and pressures in \texttt{MeV/fm$^3$}, and the evolution
step is handled in \texttt{fm} via \texttt{dtconv=dt}.

Important practical note:
unit handling is internally consistent only if the full cycle
\texttt{readinput -> phase1 -> phase2 -> phase3} is respected, because some
variables are transformed between intensive and extensive forms during
intermediate phases.

\section{Input Files}
\subsection{\texttt{input/inhydro}}
The code reads \texttt{inhydro} at startup. In the current mode, relevant lines are:
\begin{longtable}{|p{2cm}|p{11cm}|}
\hline
Line & Meaning \\\hline
1 & \texttt{ioptos1 ioptos2 ioptos3 ioptos4 kdebug kslab} \\\hline
2 & run identifier string \\\hline
3 & output file name \\\hline
4 & initial state tag/file name string (legacy field, still read) \\\hline
5 & \texttt{keqost icysum krandm tnc} \\\hline
6 & \texttt{jdim1} (geometry dimensional flag) \\\hline
7 & for \texttt{kslab>3}: \texttt{porton dtmod ncycpt ncycrp ncycaf} \\\hline
8 & \texttt{npart npartz} (read in \texttt{readinput}) \\\hline
\end{longtable}

Current repository example:
\begin{verbatim}
0 0 0 1 0 8
test for new IS
out-test
newIS-3-3-5-4
-3 301 1 10
3
.95 0.5 10 10 1
3 3
\end{verbatim}

\subsection{\texttt{input/fort.100}}
The event-by-event initialization is read from \texttt{../input/fort.100}
inside \texttt{readinput}. The file format is:
\begin{enumerate}
\item global header values (one per line):
\texttt{epsilon00, b0, ap, at, rap00, dxv, dyv, dzv};
\item then one line per active cell:
\texttt{i j k e an p w u v zmom xmom ymom}.
\end{enumerate}

The code converts GeV-based input fields to MeV-based internal fields
for the evolution pipeline.

\section{High-Level Execution Flow}
\subsection{Initialization phase}
\begin{enumerate}
\item read global control parameters from \texttt{inhydro};
\item enforce \texttt{kslab=8};
\item initialize nuclear constants and EoS tables;
\item enter \texttt{picrel}.
\end{enumerate}

\subsection{Within \texttt{picrel}}
\begin{enumerate}
\item set grid/time variables and global counters;
\item call \texttt{readinput} for event initialization;
\item build pre-evolution field snapshots;
\item run time loop up to \texttt{icysum} cycles;
\item per cycle, call:
  \begin{itemize}
  \item \texttt{phase1}: Eulerian update of cell fields;
  \item \texttt{phase2}: Lagrangian marker advection and transfer;
  \item \texttt{phase3}: reconstruction of local thermodynamic variables.
  \end{itemize}
\item build freeze-out surface with Cornelius and write hypersurface output.
\end{enumerate}

\section{Detailed PIC Workflow}
\subsection{Step A: field and marker initialization (\texttt{readinput})}
\begin{itemize}
\item The mesh dimensions \texttt{ibar, jbar, kbar} are set from memory limits.
\item Cell arrays are zeroed and then filled from \texttt{fort.100}.
\item The imported profile is shifted to the hydro mesh center.
\item Marker weight \texttt{dpt} is computed from total baryon content and
      target marker occupancy.
\item For each occupied cell, a number \texttt{mz} of markers is assigned.
\item Marker positions are stored in \texttt{px,py,pz}; cell baryon content
      is deposited into \texttt{capn} through repeated increments by \texttt{dpt}.
\end{itemize}

\subsection{Step B: local Eulerian update (\texttt{phase1})}
\begin{itemize}
\item Pressure-gradient-like contributions are evaluated from neighbor cells.
\item Momentum and energy fields are advanced using \texttt{dtodx, dtody, dtodz}.
\item Causality checks are applied (for example, handling of states with
      excessively large momentum-to-energy ratios).
\item At phase end, selected quantities are transformed to extensive form
      (multiplication by cell volume) for transfer/advection stages.
\end{itemize}

\subsection{Step C: marker transport (\texttt{phase2})}
For each marker:
\begin{enumerate}
\item read current position \texttt{(xte,yte,zte)};
\item compute effective local velocity by interpolation with \texttt{wtnwt};
\item advance marker position:
\begin{align*}
x_{n+1} &= x_n + u_{\text{eff}}\,dt/dx,\\
y_{n+1} &= y_n + v_{\text{eff}}\,dt/dy,\\
z_{n+1} &= z_n + w_{\text{eff}}\,dt/dz;
\end{align*}
\item if marker changes cell, move extensive quantities from old to new
      location using \texttt{strmom};
\item if marker exits boundaries, register as leaving matter.
\end{enumerate}

\subsection{Step D: thermodynamic reconstruction (\texttt{phase3})}
\begin{itemize}
\item Gather transferred extensive quantities per cell.
\item Convert back to intensive form by dividing by cell volume.
\item Classify interior/exterior cells with \texttt{koccupy}.
\item Reconstruct local rest-frame quantities and pressure through
      \texttt{intern}/\texttt{interx} and \texttt{eqsqcd3}.
\item Update flow velocities \texttt{u,v,w} and thermodynamic diagnostics
      (temperature, chemical potential, entropy density).
\end{itemize}

\section{Interpolation and Occupancy Logic}
\subsection{\texttt{wtnwt}}
This routine computes weighted local averages around a marker position.
It supports different weighting orders through \texttt{npower/interv}.
The output effective quantities are:
\texttt{ueff, veff, weff, eneff}.

\subsection{\texttt{koccupy}}
This routine determines if a cell is interior or exterior based on occupied
neighbors. The classification affects how reconstruction is performed in
\texttt{intern}/\texttt{interx}.

\section{Equation of State (EoS)}
\begin{itemize}
\item \texttt{read\_chiral\_eos} loads chiral-EoS tables.
\item \texttt{eqsqcd3} is used repeatedly to obtain pressure and thermal
      quantities from energy and baryon density.
\item The code contains legacy comments for bag-model forms; in the current
      branch the active path uses the tabulated/chiral EoS calls.
\end{itemize}

\section{Freeze-out Hypersurface and Cornelius}
After each full cycle, pre/post fields are available in
\texttt{/prevent/} and \texttt{/posevent/}. The code:
\begin{enumerate}
\item builds local 4D cubes with \texttt{fillcube};
\item calls Cornelius for iso-energy surface extraction (threshold \texttt{ecrit});
\item interpolates hydrodynamic quantities on each surface element;
\item writes output suitable for hadronic sampling (for example SMASH format).
\end{enumerate}

The cumulative baryon and energy flows through the surface are tracked by
\texttt{totbarcornelius} and \texttt{totenercornelius}.

\section{Output Files (Typical)}
The routine \texttt{opnfil} opens several fixed-unit files, including:
\begin{itemize}
\item \texttt{fort.6}, \texttt{fort.333}, \texttt{fort.334}, \texttt{fort.336},
      \texttt{fort.337}, \texttt{fort.338}, \texttt{fort.339}, \texttt{fort.889};
\item \texttt{oscfil.dat} for OSCAR-style output;
\item \texttt{fort\_118.dat} for cycle diagnostics.
\end{itemize}

Many outputs are controlled by \texttt{ioptos1..ioptos4}, \texttt{kdebug},
and cycle frequency parameters \texttt{ncycpt, ncycrp, ncycaf}.

\section{Build and Run Notes}
Repository \texttt{README.md} gives a baseline compilation command using
\texttt{gfortran} and \texttt{cornelius.f90}. Adapt the source filename
(\texttt{hydro\_1D\_IS.f} vs \texttt{hydro\_3D\_IS.f}) to the active branch.

A minimal run checklist is:
\begin{enumerate}
\item place \texttt{inhydro} in \texttt{input/};
\item provide \texttt{fort.100} in \texttt{input/};
\item ensure \texttt{kslab=8} in \texttt{inhydro};
\item run executable from a directory where relative paths
      \texttt{../input/*} are valid.
\end{enumerate}

\section{Maintenance Notes and Caveats}
\begin{itemize}
\item The code uses very large static arrays and extensive \texttt{COMMON} state.
\item Several comments indicate historical unit transitions; changing units in one
      phase without updating others can break conservation checks.
\item Internal/external cell logic and coasting-cell handling are tightly coupled
      to numerical stability.
\item For safe refactors, preserve the order:
\texttt{phase1 -> phase2 -> phase3}, and validate global conservation at each
stage.
\end{itemize}

\section{Subroutine Map (Quick Reference)}
\begin{longtable}{|p{3.2cm}|p{10cm}|}
\hline
Routine & Role \\\hline
\texttt{hydro} & Program entry point, input parsing, setup, call to \texttt{picrel}. \\\hline
\texttt{picrel} & Main evolution driver and cycle loop. \\\hline
\texttt{readinput} & Event initialization from \texttt{fort.100} and marker generation. \\\hline
\texttt{phase1} & Cell-field Eulerian update in conservative variables. \\\hline
\texttt{phase2} & Marker advection and cell-to-cell transfer of extensive quantities. \\\hline
\texttt{phase3} & Reconstruction of local rest-frame thermodynamics and velocities. \\\hline
\texttt{wtnwt} & Weighted interpolation of local field values around marker positions. \\\hline
\texttt{koccupy} & Interior/exterior cell classification from neighborhood occupancy. \\\hline
\texttt{intern}/\texttt{interx} & EoS-based inversion/reconstruction for exterior/interior contexts. \\\hline
\texttt{eqsqcd3} & EoS evaluator (pressure, temperature, chemical potential, entropy). \\\hline
\texttt{fillcube} & Build local 4D blocks for hypersurface extraction. \\\hline
\texttt{oschyper} & Write hypersurface elements and hydro fields to output format. \\\hline
\hline
\end{longtable}

\section{Conclusion}
The code combines an Eulerian hydrodynamic update and a Lagrangian marker
transport strategy to evolve QGP initial conditions event by event.
The implementation is highly optimized around static memory and legacy
Fortran conventions, but the computational pipeline is clear:
\begin{center}
\texttt{initialize -> phase1 -> phase2 -> phase3 -> freeze-out -> repeat}
\end{center}

This document is intended as a working baseline for new developers and for
future technical refactors.

\end{document}
